\section{Fuentes de señal de cámaras}

El sistema acepta hasta cuatro fuentes de cámara simultáneas (cam1--cam4), que
constituyen las entradas principales del mezclador. Cada cámara envía su señal al
núcleo mediante una conexión TCP en formato Matroska con vídeo en bruto (I420) y audio
PCM S16LE a 48~kHz estéreo.

\subsection{Conexión de las fuentes}

Cada fuente de cámara se conecta al puerto TCP correspondiente de voctocore:

\begin{itemize}
  \item \textbf{cam1} $\rightarrow$ puerto 10000
  \item \textbf{cam2} $\rightarrow$ puerto 10001
  \item \textbf{cam3} $\rightarrow$ puerto 10002
  \item \textbf{cam4} $\rightarrow$ puerto 10003
\end{itemize}

El envío de la señal puede realizarse con cualquier herramienta capaz de producir un
flujo Matroska sobre TCP. En el entorno de desarrollo y pruebas, las fuentes de cámara
se simulan mediante scripts FFmpeg que generan vídeo de prueba (barras de color o
\textit{test card} con reloj en tiempo real). En un despliegue real con cámaras físicas,
la señal se captura habitualmente mediante tarjetas de captura y se reencapsula en
Matroska/TCP con FFmpeg.

Los scripts de prueba se encuentran en el directorio
\texttt{example-scripts/ffmpeg/source-testvideo-as-cam*.sh}. En el escenario Docker,
cada cámara es un contenedor independiente que comparte el \textit{network namespace}
de voctocore (directiva \texttt{network\_mode: service:voctocore}), lo que les permite
conectarse directamente por \texttt{localhost} sin overhead de red virtual.

\subsection{Previsualización en la interfaz gráfica}

Una vez conectadas, las fuentes de cámara aparecen automáticamente en el panel de
previsualizaciones de voctogui. La interfaz muestra cuatro miniaturas JPEG de las
fuentes activas, recibidas desde los puertos 14000--14003 a una resolución reducida
(1024×576) para minimizar el consumo de ancho de banda entre el servidor y el cliente
de control. Cada miniatura incluye:

\begin{itemize}
  \item Un indicador visual del nombre de la fuente (cam1, cam2, cam3, cam4).
  \item Botones de selección para asignar la fuente al bus A o al bus B del mezclador.
  \item Un indicador del nivel de audio asociado a esa fuente.
\end{itemize}

El realizador selecciona la fuente activa haciendo clic sobre la miniatura
correspondiente, lo que envía el comando \texttt{set\_video\_a <fuente>} o
\texttt{set\_video\_b <fuente>} al núcleo mediante el protocolo de control TCP.

% --- Figura sugerida ---
% \begin{figure}[H]
%   \centering
%   \includegraphics[width=0.8\textwidth]{figuras/panel_camaras.png}
%   \caption{Panel de previsualizaciones de las cuatro fuentes de cámara en voctogui.}
%   \label{fig:panel_camaras}
% \end{figure}
